%----------------------------------------------------------------------
\section{Spatial Fusion: MIMD via Block-Diagonal Masked Matrix}
\label{sec:fusion}\footnote{Section-level Toulmin. \emph{Move}: introduce the first compound production, which packs heterogeneous atoms into one kernel by embedding them as blocks of a single matrix applied to a wide SIMT register. \emph{[DAF] comparison}: [DAF]'s wedge product is a bilinear operation on pairs of discriminators; the block-diagonal fusion here is the analogous construction on pairs of atoms at the scheduling level. \emph{Residue instance}: the fusion's residue is the difference between $n$ kernel launches' costs and one fused kernel's costs --- always non-negative under a monotone profile.}
%----------------------------------------------------------------------

\subsection{The block-diagonal fusion principle}

Suppose atoms $A_1, \ldots, A_n$ act on disjoint lane
ranges of a shared register. Over $\GF(2)$, their
concurrent application is a single matrix-vector
multiply:
\begin{equation}\label{eq:block-diag}
  \mathrm{reg}' = M \cdot \mathrm{reg},
  \qquad
  M = \mathrm{blockdiag}(M_{A_1}, M_{A_2}, \ldots,
  M_{A_n})
\end{equation}
where each $M_{A_i}$ is the $\GF(2)$ matrix
representation of $A_i$. The block-diagonal structure
implies that no lane of output $A_i$ depends on any
input lane outside $A_i$'s range --- they are
independent subproblems stacked into one dispatch.

\begin{definition}[\texttt{fused\_matrix\_apply} production]
  \label{def:fused-apply}
  The \emph{fused-matrix-apply} production combines
  atoms $A_1, \ldots, A_n$ with lane-disjoint masks
  $m_1, \ldots, m_n$ (where $m_i \wedge m_j = 0$ for
  $i \neq j$ in $\GF(2)$) into a single plan whose
  dispatch is one kernel applying
  \eqref{eq:block-diag}. The production's attribute
  rule is
  \begin{align}
    E &= \text{common output encoding of all } A_i \\
    \VRAM &= \max_i \VRAM(A_i) + \VRAM_{\mathrm{I/O}} \\
    \SMEM &= \sum_i \SMEM(A_i) + \SMEM_{\mathrm{mask}} \\
    \REG  &= \sum_i \REG(A_i)  \\
    \LAT  &= \max_i \LAT(A_i)
  \end{align}
  where $\VRAM_{\mathrm{I/O}}$ is the cost of loading
  the packed input register and writing the packed
  output, and $\SMEM_{\mathrm{mask}}$ is the storage
  for the masks $m_i$.
\end{definition}

\begin{proposition}[Fused-apply correctness]
  \label{prop:fused-correctness}
  \footnote{Class \textsf{AHF}. Warrant: linearity and
  block-diagonal decomposition. Qualifier: requires
  all $A_i$ to be $\GF(2)$-linear on their respective
  lane ranges; masks must be pairwise disjoint.
  Finite-$\kappa$ valid. Novelty: the application to
  compiled CSTZ kernels is new; block-diagonal
  fusion in BLAS is classical.}
  If $A_1, \ldots, A_n$ are $\GF(2)$-linear atoms
  with pairwise disjoint masks, then the output of
  \texttt{fused\_matrix\_apply}$[A_1, \ldots, A_n]$ on
  a packed input $\mathrm{reg}$ agrees, on each lane
  range, with the output of $A_i$ on the projection
  of $\mathrm{reg}$ to that range.
\end{proposition}

\begin{proof}
  By \eqref{eq:block-diag}, the output of lane $\ell$
  depends only on $M[\ell, :]$ and $\mathrm{reg}$. For
  $\ell$ in $A_i$'s range, $M[\ell, :]$ has support
  only on $A_i$'s lane range (by block-diagonal
  construction), so the output equals the output of
  $A_i$ on $\mathrm{reg}|_{\text{range}(A_i)}$. For
  $\ell$ outside any range, $M[\ell, :] = 0$ and the
  output is 0 (the identity under $\oplus$).
\end{proof}

\subsection{The LCM packing strategy}

Proposition \ref{prop:lcm-density} guaranteed that
information density per hardware tile is constant
under LCM packing. The \texttt{fused\_matrix\_apply}
production is the mechanism by which the planner
realizes that density in practice: given a workload
with atoms of mixed natural tile sizes $d_1, d_2,
\ldots$, the planner selects $L = \lcm(d_1, d_2,
\ldots, M)$ as the kernel's block size and packs
enough independent atom instances into the block to
saturate every hardware tile.

\begin{corollary}[Heterogeneous tile packing]
  \label{cor:heterogeneous}
  The \texttt{fused\_matrix\_apply} production admits
  atoms of different natural tile sizes as children,
  provided the LCM block size $L$ is a multiple of
  each. Atoms of size $d$ occupy $d \times d$
  sub-blocks of the LCM block, and the block-diagonal
  structure ensures independence.
\end{corollary}

\subsection{Mask construction}

For $n$ atoms sharing a register, the planner
constructs the masks $m_1, \ldots, m_n$ as follows.
Each $m_i$ is a $\GF(2)$-vector of length equal to the
register width, with 1s in the lanes allocated to
$A_i$ and 0s elsewhere.

\begin{proposition}[Mask disjointness is sufficient]
  \label{prop:mask-disjoint}
  If $m_i \wedge m_j = 0$ (bitwise) for all $i \neq j$,
  then the fused kernel of Proposition
  \ref{prop:fused-correctness} is a correct
  implementation of the parallel application of
  $A_1, \ldots, A_n$.
\end{proposition}

\begin{proof}
  Disjointness ensures no lane is claimed by two
  atoms. The block-diagonal structure of $M$ is a
  direct consequence: $M_{A_i}$ occupies the sub-matrix
  indexed by the supports of $m_i$, and the rest of
  $M$ is zero.
\end{proof}

\subsection{Cost vector combinator}

The cost combinator for \texttt{fused\_matrix\_apply}
has a distinctive shape:

\begin{itemize}[nosep]
  \item \textbf{VRAM} is maxed: all inputs and outputs
    are loaded once and written once, not per-atom.
  \item \textbf{SMEM} is summed: each atom may hold a
    sub-matrix in shared memory; these sub-matrices
    coexist in the fused kernel.
  \item \textbf{REG} is summed: the packed register
    holds all atoms' live values concurrently.
  \item \textbf{LAT} is maxed: the fused kernel's
    latency is the critical path, not the sum.
\end{itemize}

This is the exact complement of the cost combinator for
time-sharing (Section \ref{sec:time-sharing}), where
SMEM and REG are maxed while LAT is summed. The choice
between the two productions is the planner's most
consequential decision for register-bound workloads.

\subsection{Practical limits}

\begin{proposition}[Register-budget limit]
  \label{prop:reg-budget}
  For a hardware profile with per-thread register
  budget $B_{\REG}$, the \texttt{fused\_matrix\_apply}
  production is valid only if
  $\sum_i \REG(A_i) \leq B_{\REG}$. When this bound
  is violated, the planner must either partition the
  atoms into smaller fused groups (each under the
  budget) or substitute the \texttt{time\_share\_register}
  production (Section \ref{sec:time-sharing}).
\end{proposition}

\begin{remark}
  Exceeding $B_{\REG}$ on a real GPU triggers register
  spill to local memory (a slice of L1), with a
  concomitant latency hit on the order of 3--10$\times$
  per spilled value per access. The cost model treats
  this as an implicit edge in the cost vector's
  \LAT\ component, rising sharply once the register
  budget is saturated; in practice the planner's
  Pareto dominance check eliminates plans that trigger
  spill unless the alternatives are worse on some other
  axis.
\end{remark}

\subsection{Interaction with encoding}

Fusion respects encoding boundaries. The
\texttt{fused\_matrix\_apply} production requires all
children to share a common output encoding $E$; mixing
$\Lambda$ and RM atoms in one fused kernel is permitted
only if the planner has inserted encode/decode atoms to
align them. This is a semantic constraint, not an
efficiency one: mixing encodings would produce a
block-diagonal matrix whose rows interpret different
bases, which has no well-defined reference semantics.

Section \ref{sec:parsing} shows how the chart parser
enforces this by refusing to combine items with
mismatched $E$ attributes under the
\texttt{fused\_matrix\_apply} production rule.
